\section{主理想整环上的自由模}

定理1. --- 设A是一个环，使得A的每个左理想作为A-模都是投射的（II，第231页，定义1）。则自由左A-模L的每个子模M都是与A的理想同构的模的直和。

设 \((e_{\iota})_{\iota \in I}\) 是 L 的一个基，并令 \(p_{\iota}\) 为对应于该基的坐标函数。选取 I 的一个良序（集合论，III，第 153 页），并令 \(L_{\iota}\) 表示由诸 \(e_{\lambda}\) （ \(\lambda \leqslant \iota\) ）生成的子模；令 \(M_{\iota} = M \cap L_{\iota}\) 。坐标函数 \(p_{\iota}\) 把 \(M_{\iota}\) 映到 A 的一个理想 \(\mathfrak{a}_{\iota}\) 上；由于 \(\mathfrak{a}_{\iota}\) 是投射 A-模，存在（II，第 231 页，命题 4）一个子模 \(N_{\iota}\) （ \(M_{\iota}\) 的），使得映射 \(x \mapsto p_{\iota}(x)\) （从 \(N_{\iota}\) 到 \(\mathfrak{a}_{\iota}\) ）是双射的。设 \(M_{\iota}'\) 为 L 中由诸 \(N_{\lambda}\) （ \(\lambda \leqslant \iota\) ）生成的子模；我们将证明 \(M_{\iota}' = M_{\iota}\) 对所有 \(\iota\) 成立，由此将得出 M 由族 \((N_{\iota})_{\iota \in I}\) 生成。事实上，假设 \(M_{\lambda}' = M_{\lambda}\) 对所有 \(\lambda < \iota\) 成立；则 \(p_{\iota}(x) \in \mathfrak{a}_{\iota}\) 对所有 \(x \in M_{\iota}\) 成立；因此存在 \(y \in N_{\iota}\) 使得 \(x - y\) 是有限多个元素 \(e_{\lambda}\) （ \(\lambda < \iota\) ）的线性组合；换言之， \(x - y\) 是 \(M_{\lambda}\) （对某个 \(\lambda < \iota\) ）的元素；归纳假设表明 \(x - y \in M_{\lambda}' \subset M_{\iota}'\) ，即 \(x \in M_{\iota}'\) ，从而 \(M_{\iota}' = M_{\iota}\) 。还需证明诸 \(N_{\iota}\) 之和是直和；现在假设存在一个线性关系 \(\sum_{\iota} a_{\iota} = 0\) ，其中

\(a_{\iota} \in \mathbf{N}_{\iota}\) ，其中诸 \(a_{\iota}\) （除有限多个外均为零）并非全为零。设 \(\mu\) 为最大的指标 \(\iota\) ，使得 \(a_{\iota} \neq 0\) ；由于 \(p_{\mu}(a_{\lambda}) = 0\) （对 \(\lambda < \mu\) ），我们有 \(p_{\mu}(a_{\mu}) = p_{\mu}\left(\sum a_{\iota}\right) = 0\) ，所以 \(a_{\mu} = 0\) ，这与 \(\mu\) 的选取相矛盾.

推论1. --- 若A的每个左理想都是投射的，则投射左A-模的每个子模都是投射的。

确实每个投射A-模都是自由A-模的子模（II，第231页，命题4），因此定理1适用。

推论2. --- 主理想整环上自由模的每个子模都是自由的。

这直接由定理1推出，因为主理想整环的每个理想都是自由的。

推论3. --- 主理想整环上的每个投射模都是自由的。

注记. --- 定理 1 的证明表明， \(\mathbf{A}^{(1)}\) 的每个子模都同构于一个直和 \(\oplus \mathfrak{a}_{\iota}\) ，其中每个 \(\mathfrak{a}_{\iota}\) 是 \(\mathbf{A}\) 的理想.

命题1. --- 若L是主理想整环A上的秩为n的自由模，则L的每个子模M都是秩≤n的自由模。

确实，由定理 1 的推论 2，M 是自由模，且由前面的注记或由下面的引理，其秩为 \(\leqslant n\) ：

引理 1. --- 设 L 是交换环 A 上的一个由 n 个元素生成的模，且 M 是 L 的一个自由子模；则 M 的秩 ≤ n。

首先假设 L 是自由的。设 i 表示 M 到 L 中的典范单射。由 III 第 520 页推论，同态 \(\bigwedge_{n+1}^{n+1} i: \bigwedge_{n+1}^{n+1} M \to \bigwedge_{n+1}^{n+1} L\) 是单射；由 III 第 511 页命题 6， \(\bigwedge_{n+1}^{n+1} L = \{0\}\) ，所以 \(\bigwedge_{n+1}^{n+1} M = \{0\}\) ；由此可知 M 的秩为 \(\leqslant n\) （III，第 518 页，推论 1）。现在考虑一般情形；模 L 是一个秩为 n 的自由模 \(L'\) 的商。存在 L 的一个与 M 同构的子模 \(M'\) （II，第 218 页，命题 21）。由论证的第一部分， \(M'\) 的秩为 \(\leqslant n\) ，结论随之得出。

推论. --- 设 E 是主理想整环 A 上的模，由 n 个元素生成。则 E 的任意子模 F 可由至多 n 个元素生成。

事实上，存在一个同态 \(f\) ，从 \(\mathbf{A}^n\) 映到 \(\mathbf{E}\) 上（II，第 216 页，推论 3），且 \(f^{-1}(\mathbf{F})\) （它是秩为 \(m \leqslant n\) 的自由模）由 \(n\) 个元素生成；这些元素在 \(f\) 下的像生成 \(\mathbf{F}\) .

